\documentclass[10pt,aspectratio=169]{beamer}
% Preámbulo modular para presentaciones Beamer (Modelación Estadística)
% Diseño y paleta institucional del Tecnológico de Monterrey

\documentclass[aspectratio=169,xcolor={dvipsnames,table}]{beamer}

\usepackage[utf8]{inputenc}
\usepackage[spanish,mexico]{babel}
\usepackage{amsmath,amssymb,amsfonts,mathtools}
\usepackage{graphicx}
\usepackage{listings}
\usepackage{booktabs}
\usepackage{tikz}
\usetikzlibrary{matrix,arrows,decorations.pathmorphing,shapes.geometric,positioning}

% Paleta Institucional Tecnológico de Monterrey
\definecolor{TecAzul}{HTML}{0039A6}       % Azul TEC: color primario institucional
\definecolor{TecAzulOscuro}{HTML}{1A2E51} % Azul marino oscuro
\definecolor{TecRojo}{HTML}{EC2661}       % Rojo de acento (paleta miTec)
\definecolor{TecGris}{HTML}{646464}       % Gris neutro para comentarios y subtítulos
\definecolor{TecFondoBloque}{HTML}{F4F6F9}% Fondo suave para bloques

% Configuración de Tema Beamer
\usetheme{Boadilla}
\usecolortheme[named=TecAzulOscuro]{structure}
\setbeamercolor{alerted text}{fg=TecRojo}
\setbeamercolor{palette primary}{bg=TecAzulOscuro,fg=white}
\setbeamercolor{palette secondary}{bg=TecAzul,fg=white}
\setbeamercolor{palette tertiary}{bg=TecAzulOscuro!80!black,fg=white}
\setbeamercolor{title}{fg=TecAzulOscuro,bg=TecFondoBloque}
\setbeamercolor{frametitle}{fg=TecAzulOscuro,bg=TecFondoBloque}

% Bloques personalizados elegantes
\setbeamercolor{block title}{bg=TecAzulOscuro,fg=white}
\setbeamercolor{block body}{bg=TecFondoBloque,fg=black}
\setbeamercolor{block title alerted}{bg=TecRojo,fg=white}
\setbeamercolor{block body alerted}{bg=TecRojo!8,fg=black}
\setbeamercolor{block title example}{bg=TecAzul,fg=white}
\setbeamercolor{block body example}{bg=TecAzul!8,fg=black}

% Redondear bordes y añadir sombra sutil
\setbeamertemplate{blocks}[rounded][shadow=true]
\setbeamertemplate{navigation symbols}{}

% Pie de página limpio con número de diapositiva
\setbeamertemplate{footline}{
  \leavevmode%
  \hbox{%
  \begin{beamercolorbox}[wd=.7\paperwidth,ht=2.25ex,dp=1ex,left]{author in head/foot}%
    \hspace*{2ex}\insertshortauthor~~(\insertshortinstitute) \hfill \insertshorttitle
  \end{beamercolorbox}%
  \begin{beamercolorbox}[wd=.3\paperwidth,ht=2.25ex,dp=1ex,right]{date in head/foot}%
    \insertshortdate{}\hspace*{2em}
    \insertframenumber{} / \inserttotalframenumber\hspace*{2ex}
  \end{beamercolorbox}}%
  \vskip0pt%
}

% Configuración de Listings de Python para visualización pedagógica de código
\lstset{
    language=Python,
    basicstyle=\ttfamily\footnotesize,
    keywordstyle=\color{TecAzulOscuro}\bfseries,
    stringstyle=\color{TecRojo},
    commentstyle=\color{TecGris}\itshape,
    showstringspaces=false,
    breaklines=true,
    frame=lines,
    rulecolor=\color{TecAzulOscuro!30},
    backgroundcolor=\color{TecFondoBloque},
    aboveskip=0.5em,
    belowskip=0.5em,
    literate={á}{{\'a}}1 {é}{{\'e}}1 {í}{{\'i}}1 {ó}{{\'o}}1 {ú}{{\'u}}1
             {Á}{{\'A}}1 {É}{{\'E}}1 {Í}{{\'I}}1 {Ó}{{\'O}}1 {Ú}{{\'U}}1
             {ñ}{{\~n}}1 {Ñ}{{\~N}}1 {¿}{{?`}}1 {¡}{{!`}}1
}

% Metadatos institucionales por defecto
\title[Modelación Estadística]{Modelación Estadística para Ciencia de Datos e Ingeniería}
\author[J. Castillo Colmenares]{Juliho David Castillo Colmenares}
\institute[Tec de Monterrey]{Tecnológico de Monterrey \\ Escuela de Ingeniería y Ciencias}
\date{\today}

% Comandos y macros estadísticas compartidas para las presentaciones Beamer (Metropolis)

% Conjuntos numéricos básicos
\newcommand{\R}{\mathbb{R}}
\newcommand{\N}{\mathbb{N}}
\newcommand{\Q}{\mathbb{Q}}
\newcommand{\Z}{\mathbb{Z}}
\newcommand{\set}[1]{\left\{ #1 \right\}}
\newcommand{\abs}[1]{\left|#1\right|}
\newcommand{\norm}[1]{\left\| #1 \right\|}

% Comandos estadísticos y probabilísticos del curso
\newcommand{\Var}{\operatorname{Var}}
\newcommand{\cov}{\operatorname{Cov}}
\newcommand{\corr}{\rho}
\newcommand{\comb}[2]{\begin{pmatrix} #1 \\ #2 \end{pmatrix}}
\newcommand{\s}{\sigma}
\newcommand{\particion}{\mathfrak{P}}
\newcommand{\card}[1]{\#\left( #1 \right)}
\newcommand{\seq}[3]{\set{#1}_{#2}^{#3}}
\newcommand{\E}{\mathbb{E}}
\newcommand{\Prob}{\mathbb{P}}

% Entornos pedagógicos y cajas en español académico natural
\newenvironment{discusion}{%
  \begin{alertblock}{Pregunta de Discusión}%
}{%
  \end{alertblock}%
}

\newenvironment{reflexion}{%
  \begin{alertblock}{Reflexión para el Aula}%
}{%
  \end{alertblock}%
}

\newenvironment{paradoja}{%
  \begin{alertblock}{Paradoja de Exploración}%
}{%
  \end{alertblock}%
}

\newenvironment{motivacion}{%
  \begin{exampleblock}{Motivación: Ciencia de Datos en la Práctica}%
}{%
  \end{exampleblock}%
}

\newenvironment{retoclase}{%
  \begin{block}{Reto Rápido en Equipo}%
}{%
  \end{block}%
}


\title[Errores estándar]{Sección 6.4: Errores estándar}
\subtitle{Precisión de estimadores y margen de error}
\author[J. Castillo Colmenares]{Juliho Castillo Colmenares}
\institute{Tecnológico de Monterrey}
\date{}

\begin{document}
\begin{frame}[plain]\vspace{-0.8cm}\titlepage\end{frame}

\begin{frame}{Hoja de ruta --- capítulo 6}
  \footnotesize
  \begin{itemize}\itemsep=0.08cm
    \item \textbf{06.01--06.02} Estimación puntual y por intervalo
    \item \textbf{06.04} Errores estándar \emph{(hoy)}
    \item \textbf{06.05} Proporciones y diferencias
    \item \textbf{06.07} Tamaño de muestra
  \end{itemize}
  \begin{block}{Objetivo didáctico}
    Calcular el error estándar teórico y su estimación muestral para medias,
    diferencias y proporciones.
  \end{block}
\end{frame}

\begin{frame}{Definición y propósito}
  \footnotesize
  Para cualquier estimador puntual $\hat\theta$,
  \[
    SE(\hat\theta)=\sqrt{\operatorname{Var}(\hat\theta)}.
  \]
  \begin{alertblock}{Interpretación}
    El error estándar es la desviación de la distribución muestral y constituye
    el bloque básico del margen de error.
  \end{alertblock}
\end{frame}

\begin{frame}{Media de una población}
  \footnotesize
  Para $\hat\theta=\bar X$:
  \[
    SE(\bar X)=\frac{\sigma}{\sqrt n},
    \qquad \widehat{SE}(\bar X)=\frac{s}{\sqrt n}.
  \]
  La segunda expresión sustituye la desviación poblacional desconocida por $s$.
\end{frame}

\begin{frame}{Diferencia de medias independientes}
  \footnotesize
  Para $\hat\theta=\bar X_1-\bar X_2$:
  \[
    SE=\sqrt{\frac{\sigma_1^2}{n_1}+\frac{\sigma_2^2}{n_2}},
    \quad
    \widehat{SE}=\sqrt{\frac{S_1^2}{n_1}+\frac{S_2^2}{n_2}}.
  \]
  Si se supone varianza común, puede usarse $S_p\sqrt{1/n_1+1/n_2}$.
\end{frame}

\begin{frame}{Diferencia de medias pareadas}
  \footnotesize
  Para las diferencias $D_i$ y $\hat\theta=\bar D$:
  \[
    SE(\bar D)=\frac{\sigma_D}{\sqrt n},
    \qquad \widehat{SE}(\bar D)=\frac{S_D}{\sqrt n}.
  \]
  La pareja se reduce a una sola muestra de diferencias.
\end{frame}

\begin{frame}{Proporción y diferencia de proporciones}
  \footnotesize
  \[
    \widehat{SE}(\hat p)=\sqrt{\frac{\hat p(1-\hat p)}{n}},
  \]
  \[
    \widehat{SE}(\hat p_1-\hat p_2)=
    \sqrt{\frac{\hat p_1(1-\hat p_1)}{n_1}+
    \frac{\hat p_2(1-\hat p_2)}{n_2}}.
  \]
\end{frame}

\begin{frame}{Tabla de referencia compacta}
  \footnotesize
  \begin{tabular}{lll}
    \hline
    Parámetro & Estimador & Error estándar \\
    \hline
    $\mu$ & $\bar X$ & $\sigma/\sqrt n$ \\
    $\mu_1-\mu_2$ & $\bar X_1-\bar X_2$ & $\sqrt{\sigma_1^2/n_1+\sigma_2^2/n_2}$ \\
    $\mu_D$ & $\bar D$ & $\sigma_D/\sqrt n$ \\
    $p$ & $\hat p$ & $\sqrt{p(1-p)/n}$ \\
    \hline
  \end{tabular}
\end{frame}

\begin{frame}{Puente computacional 1/4: sustitución muestral}
  \footnotesize
  Cuando $\sigma$, $\sigma_D$ o $p$ son desconocidos, se reemplazan por $s$,
  $s_D$ o $\hat p$.
  \begin{block}{Principio}
    La forma teórica identifica la variabilidad; la forma estimada permite
    calcularla con datos.
  \end{block}
\end{frame}

\begin{frame}{Puente computacional 2/4: efecto de $n$}
  \footnotesize
  En medias y proporciones aparece $1/\sqrt n$.
  \begin{itemize}\itemsep=0.12cm
    \item cuadruplicar $n$ reduce el error estándar a la mitad;
    \item aumentar $n$ mejora precisión con rendimientos decrecientes;
    \item el error no muestral no entra en esta fórmula.
  \end{itemize}
\end{frame}

\begin{frame}{Puente computacional 3/4: margen de error}
  \footnotesize
  \[
    E=(\text{valor crítico})\times SE.
  \]
  El error estándar es la parte de variabilidad; el valor crítico incorpora el
  nivel de confianza y la distribución de referencia.
\end{frame}

\begin{frame}{Puente computacional 4/4: elección de fórmula}
  \footnotesize
  \begin{tabular}{ll}
    \hline
    Diseño & Estimador base \\
    \hline
    Una media & $\bar X$ \\
    Dos grupos independientes & $\bar X_1-\bar X_2$ \\
    Datos pareados & $\bar D$ \\
    Respuesta binaria & $\hat p$ o $\hat p_1-\hat p_2$ \\
    \hline
  \end{tabular}
\end{frame}

\begin{frame}{Aplicación derivada 1/4 --- una media}
  \footnotesize
  La tabla de la nota asigna $\bar X$ a $\mu$.
  \begin{block}{Planteamiento}
    Obtener el error estándar teórico y su estimación cuando $\sigma$ es
    desconocida.
  \end{block}
\end{frame}

\begin{frame}{Aplicación derivada 1/4 --- solución}
  \footnotesize
  \[
    SE=\frac{\sigma}{\sqrt n},\qquad \widehat{SE}=\frac{s}{\sqrt n}.
  \]
  El segundo valor se usa en un intervalo o prueba basada en la muestra.
\end{frame}

\begin{frame}{Aplicación derivada 2/4 --- dos medias}
  \footnotesize
  Para dos muestras independientes, la incertidumbre de la diferencia combina
  las dos varianzas.
  \begin{block}{Planteamiento}
    Identificar los dos tamaños muestrales y las dos desviaciones.
  \end{block}
\end{frame}

\begin{frame}{Aplicación derivada 2/4 --- solución}
  \footnotesize
  \[
    \widehat{SE}=\sqrt{\frac{S_1^2}{n_1}+\frac{S_2^2}{n_2}}.
  \]
  Si la hipótesis de varianzas iguales es razonable, la nota permite sustituir
  por la forma agrupada con $S_p$.
\end{frame}

\begin{frame}{Aplicación derivada 3/4 --- datos pareados}
  \footnotesize
  La diferencia entre cada pareja produce una sola variable $D$.
  \begin{block}{Planteamiento}
    Aplicar la fórmula de una media a las diferencias observadas.
  \end{block}
\end{frame}

\begin{frame}{Aplicación derivada 3/4 --- solución}
  \footnotesize
  \[
    \widehat{SE}(\bar D)=\frac{S_D}{\sqrt n}.
  \]
  La dependencia dentro de cada pareja se absorbe en $D_i$ antes de calcular la
  precisión.
\end{frame}

\begin{frame}{Aplicación derivada 4/4 --- proporciones}
  \footnotesize
  Para una respuesta Bernoulli, la nota usa $\hat p=X/n$.
  \begin{block}{Planteamiento}
    Estimar la variabilidad de una proporción y de la diferencia entre dos
    proporciones.
  \end{block}
\end{frame}

\begin{frame}{Aplicación derivada 4/4 --- solución}
  \footnotesize
  \[
    \widehat{SE}(\hat p)=\sqrt{\frac{\hat p(1-\hat p)}{n}}.
  \]
  Para dos grupos se suman las dos varianzas estimadas dentro de la raíz.
\end{frame}

\begin{frame}{Síntesis de la sección}
  \footnotesize
  \begin{itemize}\itemsep=0.12cm
    \item El error estándar es la desviación de un estimador.
    \item Su forma depende del diseño y del parámetro.
    \item La sustitución muestral permite calcularlo en la práctica.
    \item El margen de error agrega el valor crítico.
  \end{itemize}
\end{frame}

\begin{frame}{Transición curricular}
  \footnotesize
  \begin{block}{Lo que queda establecido}
    Toda estimación de precisión comienza identificando el estimador y su
    distribución muestral.
  \end{block}
  \begin{alertblock}{Siguiente sección}
    \textbf{06.05 Proporciones y diferencias}: intervalos Wald, Wilson y
    comparaciones entre dos tasas.
  \end{alertblock}
\end{frame}
\end{document}
